% The problem statement body, shared by hidden_hand.tex and participant_ps.tex.
You will write a Python bot that plays a card-trading game against other teams' bots. Every
bot holds a hidden hand of cards, and they trade a stock whose value is set by everyone's
cards. No finance background is needed. If you can write a Python class and you know what an
average is, you can start today.

\begin{center}
\begin{tabularx}{\linewidth}{@{}lL@{}}
\toprule
\textbf{Build}   & One Python file: a \code{Bot} class with three methods \\
\textbf{Play}    & Two bots, zero-sum, 5 rounds per deal, hundreds of deals per fight \\
\textbf{Win}     & Highest score: your PnL (profit and loss) deal by deal, on a square-root scale so luck cannot decide it \\
\textbf{Allowed} & Python 3.10+, standard library only, AI tools \\
\textbf{Submit}  & Upload to the arena (link on the WhatsApp community). Results appear within minutes \\
\bottomrule
\end{tabularx}
\end{center}

\subsection{The game}

The house shuffles one ordinary 52-card deck. Each card is worth its rank: A $=1$, 2 to 10 at
face value, J $=11$, Q $=12$, K $=13$.

\begin{itemize}
  \item \textbf{You} get 5 cards. Only you see them.
  \item \textbf{Your opponent} gets 5 cards. Only they see them.
  \item \textbf{The board} is 5 more cards, face down. One is turned face up at the start of
        each of rounds 2, 3, 4 and 5. The last one stays hidden until the end.
\end{itemize}

The stock you trade is worth \textbf{the sum of all 15 cards}: your hand, their hand and the
whole board, which averages 105. At the end every card is turned over and every trade is paid
out. If you bought below the final sum or sold above it, you made money, and your opponent lost
exactly that amount.

\subsection{A round}

Every deal has five rounds. In each round a board card is turned (rounds 2 to 5), both bots
choose a move at the same time, one bot quotes a price and the other trades on it, and then the
roles swap.

\subsubsection{Moves}

\begin{center}
\begin{tabularx}{\linewidth}{@{}l c L@{}}
\toprule
\textbf{Move} & \textbf{Cost} & \textbf{What it does} \\
\midrule
nothing & 0 & -- \\
\code{"PEEK"} & 3 & You privately see one random card from the opponent's hand. They learn
  that you peeked, but not which card. \\
\code{("SWAP", i)} & 0 & Throw away card \code{i} of your hand face up and draw a new card that
  only you see. Once per deal. \\
\bottomrule
\end{tabularx}
\end{center}

A swap changes the stock's value. If you hold 4 lots and swap a 2 for a Q, everything you hold
is worth 10 more per lot. Your opponent sees the 2 go but not what replaced it.

\subsubsection{Trading}

The \emph{maker} names a single price. The \emph{taker} may buy up to 2 lots at price $+\,2$,
sell up to 2 lots at price $-\,2$, or pass. Every round each bot quotes once and trades once.
Nobody may hold more than 10 lots long or short.

\textbf{The maker's price is a clue.} A maker usually quotes what it thinks the stock is worth,
and that depends on its hand. The taker's choice to buy or sell is a clue too. Reading these
clues is how you learn the hidden hand.

\textbf{Fairness.} Every deal is played twice with the hands swapped. Two copies of the same
bot always score exactly zero.

\subsection{A real deal}

House bot \code{sharp} against house bot \code{reader}. Reproduce it with
\code{python backtester.py sharp reader -{}-deals 1 -{}-show}.

\begin{center}
\begin{tabular}{@{}ll@{}}
\toprule
\code{sharp} holds  & \cd{Q}\ \cc{4}\ \cd{5}\ \cs{10}\ \ch{6} \quad (sum 37) \\
\code{reader} holds & \cs{7}\ \ch{4}\ \cc{3}\ \cc{7}\ \ch{9} \quad (sum 30) \\
Board, face down    & \cc{A}\ \cc{10}\ \ch{J}\ \cc{J}\ \cc{8} \\
\bottomrule
\end{tabular}
\end{center}

\begin{description}[leftmargin=0pt]
  \item[Round 1.] \code{sharp} reckons $37 + 10 \times 6.96 = 106.6$ (ten unseen cards at the
    average unseen rank) and quotes that. \code{reader} thinks the stock is cheaper and sells 2
    at 104.57. \code{reader} then quotes 101.04, and \code{sharp} buys 2 at 103.04.
  \item[Round 2.] \cc{A} turns. \code{sharp} is long, so it swaps its \cc{4} and draws a
    \cd{7}: $+3$ on everything it holds. It keeps buying from \code{reader}.
  \item[Rounds 3--5.] Prices close in and trading stops. \code{sharp} finishes long 10.
  \item[Settlement.] $40 + 30 + 41 = 111$. \code{sharp} bought at about 103 and makes $+82$.
    The same deal is then played again with the hands swapped.
\end{description}

\subsection{Your bot}

\begin{lstlisting}[language=Python]
class Bot:
    def move(self, obs):            # once per round, before trading
        return None                 # or "PEEK", or ("SWAP", index)

    def quote(self, obs):           # you are the maker
        return 104.5                # they buy at 106.5 or sell at 102.5

    def respond(self, obs, price):  # they quoted `price`
        return 2                    # +n buy n, -n sell n, 0 pass
\end{lstlisting}

A fresh \code{Bot()} is built for every deal. \code{obs} holds everything you are allowed to
know:

\begin{center}
\begin{tabularx}{\linewidth}{@{}lL@{}}
\toprule
\textbf{Field} & \textbf{Meaning} \\
\midrule
\code{round} & 1 to 5 \\
\code{hand} & your 5 cards now, as ranks 1 to 13 \\
\code{board} & board cards turned so far \\
\code{opp\_known} & opponent cards you peeked that they still hold \\
\code{unseen} & \code{unseen[r-1]} $=$ copies of rank \code{r} you have not seen anywhere
  (card counting, done for you) \\
\code{position}, \code{cash}, \code{spent} & lots held ($+$ long, $-$ short), money from
  trades, paid for moves \\
\code{tape} & every quote so far: \code{(round, "me"/"opp", price, qty)} \\
\code{discards} & every swapped-out card: \code{(round, "me"/"opp", rank)} \\
\code{peeks} & every peek: \code{(round, "me"/"opp")} \\
\bottomrule
\end{tabularx}
\end{center}

A crash, an unusable return value or a call over 50\,ms counts as doing nothing for that
action. If your \code{quote()} fails, the house quotes 105 for you.

\subsection{House bots}

Three reference bots ship with the kit and play in the arena. You may copy anything from them.

\begin{center}
\begin{tabularx}{\linewidth}{@{}lL@{}}
\toprule
\textbf{Bot} & \textbf{What it does} \\
\midrule
\code{naive} & Counts cards: my cards $+$ unseen cards $\times$ average unseen rank. Ignores the
  opponent and never moves. \textbf{The starter bot is \code{naive}.} \\
\code{reader} & \code{naive}, plus it estimates the opponent's hand from their last price \\
\code{sharp} & \code{reader}, plus it swaps a low card when long (a high one when short) and
  sizes trades by edge. Never peeks \\
\bottomrule
\end{tabularx}
\end{center}

They are far from the best possible bot. Things they never try: peeking, learning from the
opponent's trades, expecting the opponent to swap, and quoting a price that hides their own
hand.

\textbf{Test properly.} One deal can swing by 60 or more. Run 1,000 deals (about two seconds)
on two or three seeds before you believe a change helped.

\subsection{The arena}

You submit and see results on the arena; its link is posted on the WhatsApp community.

\begin{enumerate}
  \item Log in with the team key the organisers give you.
  \item Upload your \code{.py} file. It is checked at once; a failed check shows why and does
        not count against your upload limit.
  \item Your newest working file then plays every other team's newest file and the house bots
        on hidden deals. The live standings rank the score (see Evaluation), and you can watch the first
        deal of any fight, round by round.
\end{enumerate}

Code is never shown to anyone; results and replays are public. You may upload once every 10
minutes, and your newest working file is your entry.

\textbf{The final.} After the deadline every entry plays every other on a different hidden set
of deals. Only the final decides the result.

\subsection{What you submit}

One \code{.py} file per team, made from \code{starter\_bot.py}, which:

\begin{itemize}
  \item starts with these three lines, filled in:
\begin{lstlisting}
# Team Name: Your Team Name
# Members: Name 1 (Roll No.), Name 2 (Roll No.), Name 3 (Roll No.)
# Mentor: Your Mentor's Name
\end{lstlisting}
  \item defines \code{Bot} with \code{move}, \code{quote} and \code{respond};
  \item imports only \code{math}, \code{random}, \code{statistics}, \code{collections},
        \code{heapq}, \code{bisect}, \code{itertools}, \code{functools}, \code{typing};
  \item passes \code{python backtester.py -{}-validate your\_bot.py};
  \item keeps nothing from one deal to the next.
\end{itemize}

\subsection{Evaluation}

You are ranked on a \textbf{score} built deal by deal, so that one lucky deal cannot decide
anything.

\begin{enumerate}
  \item Every deal is played twice with the hands swapped. Add your PnL from the two plays and
        call it $p$. The luck of the cards mostly cancels here.
  \item Take the signed square root of $p$: a deal pair won by 100 counts $+10$, one won by
        25 counts $+5$, one lost by 16 counts $-4$.
  \item Your score is the average of these over every deal pair of every final fight:
\[
  \text{score} \;=\; \frac{1}{N}\sum_{i=1}^{N} \operatorname{sign}(p_i)\,\sqrt{|p_i|}
\]
\end{enumerate}

A single huge win, say a lucky last card while you are long 10, counts for far less than its
raw PnL, while winning deal after deal counts in full. Tested over 12 sets of 500 deals, this
separates a better bot from a worse one far more reliably than raw PnL. \code{backtester.py}
prints your score, so you can tune for exactly what is ranked.

Ties are broken on PnL per deal. Every bot plays every other bot and the three house bots the
same number of deals. House bots do not take a place.

\subsection{Rules}

\begin{enumerate}
  \item Your bot is your team's own work. Discussing ideas is fine; sharing code or another
        team's tuned numbers is not. The starter and house bots are free to copy.
  \item AI tools are allowed. Any member must be able to explain any part of the file.
  \item Mentors guide; they do not build. A bot written or tuned by a mentor is disqualified.
  \item Play the game, not the system. Reading files, using the network or keeping state
        between deals gets the file rejected. Report loopholes; do not use them.
  \item One team, one entry. A member may be in only one team.
  \item The organisers' decision is final. Ask on the group before building on a guess.
\end{enumerate}

\subsection{Quick start}

Download the kit from the arena (\emph{Kit}, top bar), unzip it, then:

\begin{lstlisting}[language=bash]
cd hiddenhand                     # Python 3.10+, nothing to install
python backtester.py --show       # sharp vs reader, one deal printed
cp starter_bot.py my_bot.py       # your entry starts here
python backtester.py my_bot.py naive --deals 1000
python backtester.py --validate my_bot.py
\end{lstlisting}

\subsection{FAQ}

\begin{description}[leftmargin=0pt,style=nextline]
  \item[I have never traded. Is that a problem?] No. It is a probability game more than a
    finance one.
  \item[Can I use numpy or pandas?] No, standard library only. The game is small enough not to
    need them.
  \item[How fast must my bot be?] A millisecond or two per call. A call over 50\,ms counts as
    doing nothing.
  \item[Why did my rank change when I did not upload?] Other teams uploaded, and your bot now
    plays their new bots.
  \item[Where do I ask?] Your mentor first, then the WhatsApp community. Rule answers given
    there apply to everyone.
\end{description}

\subsubsection{Glossary}

\begin{center}
\begin{tabularx}{\linewidth}{@{}lL@{}}
\toprule
\textbf{PnL} & Profit and loss: how much you made (positive) or lost (negative) \\
\textbf{Long / short} & Long: you bought, so you gain if the sum ends higher. Short: you sold \\
\textbf{Lot} & One unit of the stock; each lot pays the final sum \\
\textbf{Maker / taker} & The maker names the price; the taker decides whether to trade \\
\textbf{Spread} & The gap between the buy price (price $+\,2$) and the sell price (price $-\,2$) \\
\textbf{Card counting} & Tracking which cards are gone, so you know what is left \\
\bottomrule
\end{tabularx}
\end{center}

% ═════════════════════════════════════════════════════════
